% ==============================================================================
% SLIDES - SEMANA 02: GOVERNANÇA CORPORATIVA, SOX, COMPLIANCE E ÉTICA NA TI
% ==============================================================================

% ------------------------------------------------------------------------------
% 1. SLIDE DE CAPA
% ------------------------------------------------------------------------------
\senaicover
  {Governança Corporativa \& SOX}
  {Princípios IBGC, Controles Internos (ITGC), Compliance e Ética na TI}
  {Unidade Curricular: Governança de TI}
  {Prof. André Souza}

% ------------------------------------------------------------------------------
% 2. VISÃO GERAL / AGENDA
% ------------------------------------------------------------------------------
\begin{senaiframe}{Visão Geral \& Agenda}{Os 4 Eixos da Aula}

  \begin{columns}[t]
    \begin{column}{0.48\textwidth}
      \begin{tcolorbox}[senaibluecard, title={1. Governança Corporativa}]
        \begin{itemize}
          \item Teoria da Agência e conflitos de interesse.
          \item Os 4 Princípios Universais do IBGC.
        \end{itemize}
      \end{tcolorbox}
      
      \vspace{4pt}
      
      \begin{tcolorbox}[senaiambercard, title={3. Compliance \& Integridade}]
        \begin{itemize}
          \item Conformidade legal (LGPD, BACEN, PCI).
          \item O custo estratégico da não-conformidade.
        \end{itemize}
      \end{tcolorbox}
    \end{column}

    \begin{column}{0.48\textwidth}
      \begin{tcolorbox}[senairedcard, title={2. Lei Sarbanes-Oxley (SOX)}]
        \begin{itemize}
          \item Seções 302/404 e responsabilização.
          \item Controles Gerais de TI (ITGC).
        \end{itemize}
      \end{tcolorbox}
      
      \vspace{4pt}
      
      \begin{tcolorbox}[senaigreencard, title={4. Capítulo 1 do PDGTI}]
        \begin{itemize}
          \item Caracterização da Organização-Alvo.
          \item Mapeamento da estrutura de TI e negócio.
        \end{itemize}
      \end{tcolorbox}
    \end{column}
  \end{columns}

\end{senaiframe}

% ------------------------------------------------------------------------------
% 3. OS 4 PRINCÍPIOS DO IBGC NA TI
% ------------------------------------------------------------------------------
\begin{senaiframe}{Módulo 1 • Diretrizes do IBGC}{Os 4 Princípios de Governança Aplicados à TI}

  \begin{columns}[t]
    \begin{column}{0.48\textwidth}
      \begin{tcolorbox}[senaibluecard, title={Transparência \& Equidade}]
        \begin{itemize}
          \item \textbf{Transparência:} Relatórios de SLA, custos de nuvem (FinOps) e status de projetos sem ocultar falhas.
          \item \textbf{Equidade:} Atendimento isonômico e priorização imparcial no portfólio de TI.
        \end{itemize}
      \end{tcolorbox}
    \end{column}

    \begin{column}{0.48\textwidth}
      \begin{tcolorbox}[senaicard, title={Accountability \& Responsabilidade}]
        \begin{itemize}
          \item \textbf{Accountability:} Donos de dados e sistemas (Data Owners) e Matriz RACI explicita.
          \item \textbf{Responsabilidade:} TI Verde, descarte de e-waste e proteção da privacidade.
        \end{itemize}
      \end{tcolorbox}
    \end{column}
  \end{columns}

\end{senaiframe}

% ------------------------------------------------------------------------------
% 4. SOX & ITGC (CONTROLES GERAIS DE TI)
% ------------------------------------------------------------------------------
\begin{senaiframe}{Módulo 2 • Lei SOX \& TI}{Os 4 Domínios de Controles Gerais de TI (ITGC)}

  \begin{columns}[t]
    \begin{column}{0.48\textwidth}
      \begin{tcolorbox}[senairedcard, title={Acessos \& Mudanças}]
        \begin{itemize}
          \item \textbf{Gestão de Acessos Lógicos:} Menor privilégio, MFA e revogação imediata no desligamento.
          \item \textbf{Gestão de Mudanças:} Segregação de ambientes; nenhum dev altera código direto em produção.
        \end{itemize}
      \end{tcolorbox}
    \end{column}

    \begin{column}{0.48\textwidth}
      \begin{tcolorbox}[senaicard, title={Operações \& Desenvolvimento}]
        \begin{itemize}
          \item \textbf{Operações de TI:} Backups diários imutáveis e testes anuais de Disaster Recovery (DRP).
          \item \textbf{Desenvolvimento:} Versionamento seguro (Git) e homologação formal antes do go-live.
        \end{itemize}
      \end{tcolorbox}
    \end{column}
  \end{columns}

\end{senaiframe}

% ------------------------------------------------------------------------------
% 5. REFERÊNCIAS BIBLIOGRÁFICAS
% ------------------------------------------------------------------------------
\begin{senaiframe}{Referências Bibliográficas}{Fontes \& Leitura Recomendada}

  \begin{tcolorbox}[senaicard, title={Obras de Referência}]
    \begin{itemize}
      \item \textbf{BITTENCOURT, Carlos Magno Andriolli.} *Governança corporativa e compliance: planejamento e gestão estratégica.* Curitiba: Contentus, 2020.
      \item \textbf{MARTINS, Camila Saldanha.} *Governança e compliance.* Curitiba: Contentus, 2020.
      \item \textbf{IBGC.} *Código das Melhores Práticas de Governança Corporativa.* 6. ed. São Paulo: IBGC, 2023.
    \end{itemize}
  \end{tcolorbox}

\end{senaiframe}

% ------------------------------------------------------------------------------
% 6. APLICAÇÃO NO PDGTI
% ------------------------------------------------------------------------------
\begin{senaiframe}{Aplicação no PDGTI}{Desenvolvimento do Capítulo 1}

  \vspace{0.2cm}
  \begin{tcolorbox}[senaigreencard, title={Capítulo 1 – Caracterização da Organização}]
    \begin{itemize}
      \item \textbf{1.1 Visão Geral do Negócio:} Setor, porte, missão, visão e valores da empresa.
      \item \textbf{1.2 Estrutura Organizacional:} Organograma corporativo e diretoria.
      \item \textbf{1.3 Estrutura de TI:} Organograma atual da TI e quantitativo de colaboradores.
      \item \textbf{1.4 Inventário Simplificado:} Sistemas principais (ERP/CRM), infraestrutura e orçamento.
    \end{itemize}
  \end{tcolorbox}

\end{senaiframe}
